\documentclass[9pt]{article}
\usepackage[a4paper,margin=0.25in]{geometry}
\usepackage{amsmath,amssymb,amsfonts,mathtools}
\usepackage{xcolor}
\allowdisplaybreaks
\setlength{\jot}{3pt}

\begin{document}
\section*{All three-$\rho$ terms proportional to $\log(\vee/\Lambda)$, with c.c.: final form}
One equation. The lines are numbered as in your list; $\mathcal D$ has the lines D1--D8.
\begin{itemize}
\item \textbf{Black:} exactly as you wrote it.
\item \textcolor{red}{\textbf{Red: removed}}, kept only for reference. These terms are \emph{not} part of the sum.
\begin{itemize}
\item[F4] Group I, Row II: the $\alpha$ part of the $U(w)\mathbb A^{(3)}$ sub-terms. It multiplies $-\log(\vee/k^+)$, not $\log(\vee/\Lambda)$.
\item[F3] Group I, Row III, part 2: the $\gamma,\gamma'$ terms with the edge kernel $\phi$.
\end{itemize}
\item \textcolor{blue}{\textbf{Blue: changed or added}}.
\begin{itemize}
\item[T1] Group III, Row I, line 15: $\rho^{h}(x')\to\rho^{c}(x')$.
\item[T2] Group III, Row II, line 13: $U^{hl}(y)\to U^{hl}(x)$.
\item[T3] Group I, Row II, part 1, $\alpha$ lines: the kernel $\mathcal K(x'-w)\cdot\mathcal K(y-z)\to\mathcal K(x-z)\cdot\mathcal K(y-z)$.
\item[F3] $\mathcal D$, lines D1--D8: added. D5--D8 are its c.c., with the charges in reverse order.
\item[c.c.] The c.c.\ of an operator product reverses the order of the charges.
\item[order] The anticommutator of the notes, $\rho\{\rho,\rho\}$ or $\{\rho,\rho\}\rho$, with half of your coefficient, in place of one ordering with the full coefficient. In the c.c.\ it is also reversed.
\end{itemize}
\end{itemize}
The final result is the black and blue lines. Their symmetrized three-$\rho$ parts cancel exactly. Reordering them gives the two-$\rho$ terms that, together with the genuine two-$\rho$ terms and those of \texttt{4rho.tex}, are JIMWLK$\otimes$LO.

The order of the charges only matters for the two-$\rho$ terms; for the three-$\rho$ cancellation any order gives the same result. In Group I, Row II the two orderings of a pair are exactly $\rho\{\rho,\rho\}$ through the antisymmetry of $f^{abc}$, so the amplitude is kept as written and only the c.c.\ is reversed.

@@EQ@@
\end{document}
